\documentclass[11pt]{article}
\usepackage[margin=0.75in]{geometry}
\usepackage{amsmath,amssymb}
\usepackage{graphicx}
\usepackage{booktabs}
\usepackage{microtype}
\usepackage[hidelinks]{hyperref}
\graphicspath{{./}}

\setlength{\parskip}{1pt}
\setlength{\abovecaptionskip}{1pt}
\setlength{\belowcaptionskip}{0pt}
\setlength{\textfloatsep}{4pt plus 2pt minus 2pt}   % trim whitespace above/below top-bottom floats
\setlength{\intextsep}{4pt plus 2pt minus 2pt}      % and around here-placed floats
\setlength{\dbltextfloatsep}{4pt plus 2pt minus 2pt}
\emergencystretch=1.5em  % absorb the few long \texttt{} terms cleanly
\newcommand{\code}[1]{\texttt{#1}}
\newcommand{\E}{\mathbb{E}}
\newcommand{\NPV}{\mathrm{NPV}}
\newcommand{\PD}{\mathrm{PD}}

\title{\vspace{-1.5em}NPV-Optimal SMB Underwriting under Selection Bias,\\
with Calibrated Risk and Interventional Counterfactuals}
\author{\textbf{Team Vector Ventures}}
\date{June 2026}

\begin{document}
\maketitle
\vspace{-1.5em}

\section{Problem framing \& assumptions violated}

\emph{Every assumption a credit scorecard leans on---random labels, population
overlap, a stationary window, honest inputs---breaks on this book; we name each break
and respond to it rather than modelling straight through it.}

We are re-underwriting a historical SMB loan book to maximize realized portfolio net
present value (NPV)---total dollars earned, not classification accuracy.

\textbf{Labels are missing-not-at-random behind a sharp gate.} Outcomes exist only for
loans the prior lender approved, and approval is, to numerical precision, a hard cutoff
$\code{prior\_underwriter\_score} \geq 0.273$: a propensity model separates approved from
declined with AUC $1.00$ and \emph{zero} score overlap. So the label is selected on a
deterministic function of an observed covariate---textbook sample-selection bias. Inverse-propensity
reweighting is infeasible here by a positivity violation, \emph{not by our choice}: a propensity
model on the ordinary covariates alone still hits AUC $1.00$, so approved and declined applicants
share no common feature support. Only local extrapolation near the cutoff is available, and we
widen intervals on declines rather than pretend we can reweight.

\textbf{The split is forward in time.} Training spans an 18-month book ending the day before
the 13-week scoring window opens, over which the default rate drifts $17.5\% \to 20.6\%$. We
learn default \emph{timing} on the history but recalibrate the \emph{level} in-window---the
history tells us the shape, only the window tells us the rate.

\textbf{Missingness is information, not absence.} ``Never declined elsewhere'' (null) loans
default at $14.3\%$ versus $21.3\%$ when present; no-bank-feed loans at $19.1\%$ versus
$16.7\%$---a blank field predicts the outcome, so we encode it as a first-class indicator
rather than impute it away. \textbf{Self-reports are optimistically biased}: applicants who
overstate revenue by more than $1.5\times$ default at $38.6\%$ versus $16.1\%$, so the
\emph{gap} between stated and bank-feed revenue---not the stated figure---is the signal,
central to our causal treatment of \code{stated\_annual\_revenue} below.

\section{Methodology}

\emph{One discrete-time model serves all three deliverables: a calibrated probability
that a loan defaults, times a per-band curve for \emph{when} it defaults, read out as an
NPV sign for the approve/decline call.}

\textbf{One shared discrete-time model.} The brief's NPV depends on the \emph{day} of
default (a day-5 default loses roughly the principal; a day-55 default nearly breaks
even), and Deliverable B is the cumulative-default curve. Both consume one timing
model. We write each loan's cumulative-default curve as
$F_i(t) = \PD_i \cdot S_{b(i)}(t)$, where $\PD_i = \Pr(y_i{=}1 \mid \mathbf{x}_i)$ is a
calibrated probability of default by day~90 and $S_{b}(t)$ is a normalized timing
shape ($S_b(90){=}1$) learned per credit band $b$. In words: \emph{how likely} a default
is, kept separate from \emph{when} it arrives.

\textbf{PD model.} A five-fold \code{GroupKFold} (by \code{business\_id}) LightGBM
ensemble over 55 engineered features (out-of-fold AUC $0.774$, Brier $0.117$),
isotonic-calibrated. Features are deliberately buffer-centric: because loans are
$\sim$1.3\% of annual revenue, default is cash-buffer bound rather than revenue bound, so
debt-service coverage, buffer-to-payment, overdrafts and credit utilization dominate.

\textbf{The timing shape is bimodal and we model it as such} (Figure~\ref{fig:traj}).
Of all defaults, $77.5\%$ are missed-draw events over days $3$--$60$; \emph{no} defaults
occur on days $61$--$89$; then the remaining $22.5\%$ spike exactly at day~90 (the
open-balance sweep). A smooth survival fit would smear that mass across weeks, so we
encode the empirical concave-rise/flat/day-90-cliff shape rather than blur the cliff. The shape is
band-conditional---worse-credit borrowers default $\sim$13 days earlier and carry less
day-90 mass---so $S_b(t)$ feeds both B and the per-loan $\E[t^*]$ used in A.

\begin{figure}[t]
\centering
\includegraphics[width=0.40\textwidth]{fig03_cumulative_trajectory.png}
\caption{Pooled cumulative default rate by loan-age week: a concave rise through
week~9, a flat stretch (weeks 9--12), and the day-90 open-balance cliff at week~13.
Deliverable~B reproduces this shape per cohort; a single smooth survival curve cannot.}
\label{fig:traj}
\end{figure}

\textbf{Deliverable A: decide on the NPV sign.} We approve iff
$\E[\NPV_i \mid \text{approve}] > 0$ using the exact product economics: revenue
$F_i + R_i r (T/365)$ if repaid, and $F_i + D_i(t^*_i - 1) + \mathrm{rec}_i - R_i$ if
default at day $t^*_i$---we keep fee$+$interest if repaid, but only fee plus pre-$t^*$ draws
(less principal, $+9\%$ recovery) on a default, so later defaults cost less. NPV is linear in
$t^*$, so the expectation integrates the band-conditional default-day distribution exactly. The
break-even sits at PD $\approx 0.39$, not $0.5$: even defaulters repay $\sim$74\% of the schedule
first, so a $39\%$-PD applicant is still NPV-positive. We then shift the decision PD up by
$\kappa\hat\sigma_i$ ($\hat\sigma_i$ the fold disagreement, $\kappa{=}1.25$ by cross-fit on realized
val P\&L: out-of-fold \$3.76M vs \$3.73M at $\kappa{=}0$), declining the most uncertain
near-break-even loans first. We approve 72\% of applicants; on realized validation the book
returns $1.76\times$ the prior underwriter's P\&L.

\textbf{Deliverable B: shape $\times$ level.}
$\mathrm{CDR}_{w,a} = \mathrm{mean}_{i \in A_w}\, \PD_i \, S_{b(i)}(a)$---we average each
approved cohort $A_w$'s per-loan risk over the timing shape, giving a curve that only rises.
Because the PD model has no cohort signal (training predates the window), we shrink each
cohort's \emph{level} toward the realized validation rate for the same calendar week
(pseudo-count 75) and its \emph{shape} toward the model band shape (Dirichlet $c{=}25$,
split-half cross-fit), which pulls the sparse tail cohorts onto the data (MAE $0.013{\to}0.004$).
Validation and test span the identical 13 weeks, so this transfers.

\section{Causal reasoning \& counterfactual methodology}

\emph{The judges' counterfactual queries ask what an intervention \emph{causes}, not what a
feature \emph{predicts}---so we classify every feature by whether editing it changes reality,
and scale the effect accordingly.}

\begin{figure}[t]
\centering
\includegraphics[width=0.34\textwidth]{shap_summary.png}
\caption{SHAP summary for the PD model: mean $|$SHAP$|$ ranking with per-feature
direction (red high / blue low). High delinquency, utilization and volatility raise PD;
high debt-service coverage and cash buffer lower it.}
\label{fig:shap}
\end{figure}

The counterfactual target is
$p_q^{\mathrm{cf}} = \Pr\!\big(y_{i_q}{=}1 \mid \mathrm{do}(f_q{=}v_q),\,
\mathbf{X}_{i_q,-f_q}{=}\mathbf{x}_{i_q,-f_q}\big)$ versus the observational
$\Pr(y \mid f{=}v, \mathbf{X}_{-f})$: the first asks what happens if we \emph{set} $f$ to $v$,
the second describes applicants who happen to \emph{have} $f{=}v$, and they diverge exactly
when $f$ is confounded. Our \textbf{feature registry} keeps the edit honest: each engineered
feature declares its parents, so setting a raw feature recomputes exactly its downstream
features and nothing else---never an internally-contradictory applicant carrying a \$20k loan
with a \$10k loan's payments. For example $\mathrm{do}(\code{requested\_amount})$ recomputes the
daily-payment, buffer-to-payment, leverage and ask-to-revenue features; the provided
\code{requested\_amount\_to\_observed\_revenue} is never reused (no stale-descendant bug).

\textbf{We classify the intervenable features by causal status.} \emph{Mechanical causes}
(\code{requested\_amount}): changing the number changes the borrower's actual payment burden,
so we perturb and propagate it through the registry at full strength. \emph{Confounded
proxies}---bureau and bank-feed symptoms of latent health (utilization, delinquency, revenue
and cash signals)---are symptoms as much as causes: editing the thermometer does not cure the
fever, so we keep only the estimated causal fraction
$\hat\lambda = \beta_{\text{adj}}/\beta_{\text{naive}}$ (sibling-adjusted logistic; per-family
$\approx 0.52/0.51/0.64$ for bureau/bank-feed/behavioral), averaged with the heuristic, falling
back where $\hat\lambda \notin (0,1)$. Conditioning sets are derived
\emph{algorithmically} from an explicit DAG by the backdoor criterion---condition on a
treatment's co-proxies of latent health, exclude its graph \emph{descendants}
(Figure~\ref{fig:dag}). The FDR partial-correlation graph flags ``siblings'' that are really
descendants---invoice delinquency and overdrafts follow low cash ($-0.78$), payroll regularity
follows revenue ($+0.50$); PC and GES corroborate. Since conditioning on a descendant blocks the
true causal path, excluding it by graph reachability corrects $\hat\lambda_{\text{cash}}$ from an
implausible $0.03$ to $0.44$.
\emph{Self-reports}---\code{stated\_annual\_revenue}, \code{stated\_time\_in\_business}---are
returned at their observational PD: a different number on the form changes nothing about the
business. Overstating revenue \emph{predicts} default ($38.6\%$ vs $16.1\%$), but \emph{forcing}
a different statement causes nothing, so the interventional effect is $\approx 0$---formally
$\Pr(y \mid \mathrm{do}(\text{stated}), \mathbf{X}_{-\text{stated}}) =
\Pr(y \mid \mathbf{X}_{-\text{stated}})$, the causally correct answer, not naive re-prediction.
\emph{Immutable/proxy-by-fiat} features (\code{sector}, \code{geography\_region},
\code{vintage\_years}, \code{has\_linked\_bank\_feed}) are not levers but are queried anyway; we
perturb the model, shrink, widen, and flag them.

\begin{figure}[t]
\centering
\includegraphics[width=0.92\textwidth]{causal_dag.png}
\vspace{-0.6em}
\caption{The explicit DAG. Latent health $H$ confounds every proxy (faint), so we adjust for a
treatment's co-proxies of $H$; the three defended mechanistic edges (red) make invoice
delinquency/overdrafts descendants of cash and payroll a descendant of revenue, which graph
reachability excludes from those treatments' adjustment sets. The FDR partial-correlation graph
($q{=}0.05$, $-0.78$/$+0.50$ edges) and PC/GES corroborate this structure.}
\label{fig:dag}
\end{figure}

\textbf{Regulator defense.} Drivers are explained with SHAP over an inspectable,
de-correlated feature set (Figure~\ref{fig:shap}). The leading drivers---invoice-payment
delinquency, credit utilization, revenue volatility, requested amount, overdrafts, and the
affordability ratios---are all recognized credit/affordability signals in the expected
direction, none a prohibited or proxy-discriminatory attribute. Several queried
values lie outside observed support (e.g.\ \code{prior\_loans\_count} is in-support only 29\%
of the time); these receive wider intervals and an explicit extrapolation caveat.

\textbf{Model-class check.} A transparent logistic regression on the same engineered
features matches the gradient-boosted models (out-of-fold AUC $0.775$ vs $0.770$, equal
Brier) and is in fact \emph{more} robust to the train-to-window drift (validation AUC
$0.759$ vs $0.751$)---evidence the engineered space is near-linearly sufficient and the
drivers directly interpretable. Its coefficients agree in sign with SHAP on 8/10 top features
(the two flips, \code{buffer\_to\_payment} and \code{debt\_to\_revenue}, are
collinearity-driven---why we read SHAP, not raw coefficients). We rejected sequence/transformer models: no sequential structure,
low-cardinality categoricals, $\approx\!0.77$-AUC data-limited discrimination, and calibration
favours the simpler stack.

\section{Calibration \& uncertainty quantification}

\emph{We fix the probability \emph{level} on the only data drawn from the scoring window, then
size every interval by measured out-of-sample error rather than in-sample optimism.}

\textbf{Point PD.} Isotonic calibration is \emph{refit on in-window validation} to correct the
$17.5\%{\to}20.6\%$ train-to-window drift (power-prior $a_0{\to}0$ limit); it stays valid on the
approved manifold and is extrapolated on declines.

\textbf{Intervals (A, C).} Additive on the calibrated point, $p \pm \alpha z \sigma$: the point
probability widened by how much the five folds disagree ($\sigma$), with the multiplier $\alpha$
fit by five-fold \emph{cross-fit} on validation. Validated decile coverage is $\approx 0.90$ at mean width $\approx 0.06$---the brief's
``contain the truth without being needlessly wide.''

\textbf{Trajectory intervals (B).} A within-cohort bootstrap combined with a conformal band
calibrated against the \emph{true} validation cohort trajectories; validated coverage
$\approx 0.92$. Out-of-support counterfactual intervals in C are widened further
(Table~\ref{tab:val}).

\begin{table}[h]
\centering
\begin{tabular}{@{}lll@{}}
\toprule
Deliverable & Metric & Value \\
\midrule
A --- PD model & out-of-fold AUC / Brier & $0.774$ / $0.117$ \\
A --- policy & approve rate / P\&L vs.\ prior underwriter & $0.72$ / $1.76\times$ \\
A --- intervals & decile coverage / mean width & $0.90$ / $0.06$ \\
B --- trajectory & mean abs.\ error vs.\ realized CDR / coverage & $0.004$ / $0.92$ \\
\bottomrule
\end{tabular}
\caption{Validation evidence (in-window). P\&L and B error use realized validation outcomes;
interval coverage is out-of-sample (cross-fit / conformal).}
\label{tab:val}
\end{table}

\section{Limitations \& what we'd do differently}

\emph{Our honest gaps trace to one fact---we never observe outcomes for declined
applicants---so the decline and counterfactual estimates rest on extrapolation we cannot check.}

\textbf{The unverifiable core.} All C effects lack labels (the queried applicants were never
decided), and declines never overlap the approved manifold (\S1)---so reject inference is
infeasible, only RD-style local extrapolation remains (deferred), and decline-side PD, part of
A's score, stays unchecked. Self-report and immutable counterfactuals stay heuristic; confounded
proxies use the DAG-derived $\hat\lambda$ above, but latent $H$ breaks causal sufficiency---PC/GES
corroborate the structure, they do not prove it.

\looseness=-1
\textbf{Our val-tuned knobs} ($\kappa{=}1.25$, shape $c{=}25$, per-cohort level/shape) transfer to
test only as far as the windows match (same 13 weeks; per-cohort $n\approx150$ adds variance). The
$\kappa$ rule absorbs the flagged over-approval, but a residual oracle gap remains ($\approx 0.58$
of perfect-foresight P\&L), irreducible without label foresight on declines.

\looseness=-1
\textbf{With another day:} RD extrapolation at the $0.273$ cutoff into declines, and one
discrete-time hazard unifying PD and timing.

\end{document}
